\CWHeader{Лабораторная работа \textnumero 2}

\CWProblem{
Требуется разработать программную библиотеку, реализующую словарь на основе 
структуры данных PATRICIA (Practical Algorithm To Retrieve Information Coded In Alphanumeric), 
а также программу для работы с данным словарём.

Словарь хранит пары <<ключ-значение>>, где ключом является регистронезависимая 
строка, состоящая из букв английского алфавита длиной не более 256 символов, 
а значением — целое число в диапазоне от 0 до $2^{64} - 1$. Разным ключам может 
соответствовать одинаковое значение.

Программа должна обрабатывать входной поток построчно до конца файла. Поддерживаются 
следующие операции:

\begin{itemize}
    \item \texttt{+ word value} — добавить слово в словарь. \\
    Вывод: \texttt{OK}, если вставка успешна; \texttt{Exist}, если ключ уже существует.
    
    \item \texttt{- word} — удалить слово из словаря. \\
    Вывод: \texttt{OK}, если слово было удалено; \texttt{NoSuchWord}, если слово не найдено.
    
    \item \texttt{word} — поиск слова в словаре. \\
    Вывод: \texttt{OK: value}, если слово найдено; \texttt{NoSuchWord}, если отсутствует.
    
    \item \texttt{! Save path} — сохранить словарь в бинарном виде в файл. \\
    Вывод: \texttt{OK} при успехе или \texttt{ERROR: <message>} при ошибке.
    
    \item \texttt{! Load path} — загрузить словарь из файла. \\
    При успешной загрузке текущий словарь заменяется. \\
    Вывод: \texttt{OK} или \texttt{ERROR: <message>} при ошибке.
\end{itemize}

Во всех случаях возникновения системных ошибок программа должна выводить сообщение 
в формате \texttt{ERROR: <message>}.

{\bfseries Вариант структуры данных:} PATRICIA (битовое сжатое префиксное дерево).

{\bfseries Ограничения:}
\begin{itemize}
    \item Время работы: не более 15 секунд
    \item Память: не более 512 Мб
    \item Использование стандартных ассоциативных контейнеров запрещено
\end{itemize}
}